\section{Discussion}

\subsection{Screening and foreclosure are distinct}

The model resolves the puzzle posed in the Introduction by separating two consequences that strictness can bundle in observed organizations. Club-goods theory explains why demanding practices can strengthen a group through screening: costly requirements select participants willing to invest, deter free riding, and redirect members' activity toward collective production. Those same requirements can also reduce the value or availability of life outside the group. The first process concerns who participates and whether an enforcement apparatus can operate. The second concerns whether participants who wish to leave retain a usable outside option. They can coexist in the same institution, but they are not the same causal mechanism.

In the model, activation of an enforcement apparatus is governed by code observability and the reward attached to enforcement. Retention is governed by exit capacity. Durable control requires both activation and retention. The necessity result makes the first condition visible: the low-observability, low-reward cell remains at zero capture at every imposed closure level, including the maximum. Foreclosing exit cannot manufacture an enforcement apparatus where conduct is insufficiently legible and enforcement is insufficiently rewarded. The converse result makes the second condition equally clear. With exit open, 0 of 1,890 runs reach capture, even though the overall active rate is 0.598. An apparatus can therefore activate and punish without retaining its subjects.

The implication is narrower and more useful than a general association between strictness and durability. Strictness that screens does not by itself produce lock-in. Screening can assemble committed participants and support active enforcement, but lock-in requires exit to be independently foreclosed. Conversely, closure can retain members without producing an active apparatus in structurally unfavorable cells. The two dimensions are largely separable, although they are coupled at the frontier margin, as discussed below. Durable control appears only when the organizational conditions for activation coincide with structural conditions that prevent departure.

\subsection{What the model does not show}

First, capture does not arise endogenously in this model. The earlier coupling made outside-option degradation a function of enforcer share, while enforcer share was also a component of the capture criterion. That construction was circular. With this coupling removed, capture is zero at every tested coupling strength. The model therefore does not establish a spontaneous transition from ordinary organizational activity to closed, enforcement-dominant control. Its positive results concern what happens when exit capacity is imposed as a structural condition, not how that condition originates from within the modeled interactions.

Second, punishment concentration is not emergent in the reported architecture. Roughly 84 percent of top-five-percent concentration is manufactured by the privilege bundle. Exclusive or asymmetric sanctioning opportunity, reduced costs and risks, and material reinforcement direct punishment toward a designated cadre by construction. Concentration also fails to distinguish captured systems from systems that are merely active. The top-five-percent share is 0.678 in captured systems and 0.625 in merely active systems, while the corresponding Gini values are 0.918 and 0.903. By contrast, exit rate differs sharply, at 0.105 versus 0.339. [VERIFY: the CAPTURE-versus-MIXED top-five-percent shares, Gini values, and exit rates are supplied for this draft but are not yet recorded in RESULTS_TABLE.md.] Concentrated punishment is therefore evidence about institutional allocation of coercive opportunity, not evidence that capture has occurred.

Third, the paper makes no claim that codified belief tends toward coercion. Codification can support coordination, shared identity, instruction, or collective production without producing either an enforcement apparatus or blocked exit. The mechanism is conditional and structural. It identifies a conjunction of observability, enforcement reward, institutional privilege, and exit foreclosure under which durable control can occur. It does not assign a coercive tendency to doctrine itself, and it does not infer organizational outcomes from the content or truth of belief.

\subsection{Interpretation}

Exit foreclosure should not be interpreted as an emanation of doctrine. In the framework, it is produced by identifiable conditions outside doctrinal content: legal prohibition of departure, economic dependence on the group, severance of outside ties, and generational replacement through which members are raised without outside relationships. These conditions reduce the feasibility of departure by altering opportunities and resources rather than simply changing preferences. They also clarify why exogenous treatment is analytically useful. Exit capacity can be scored as a feature of the institutional and social environment before examining whether enforcement becomes durable.

The relevant conditions are, in principle, observable in the historical record independently of the outcome. Legal rules can be distinguished from subsequent compliance, economic dependence from continued membership, network severance from later retention, and childhood socialization from adult enforcement activity. This separation makes the framework testable rather than interpretive. A study need not infer foreclosure from the persistence it seeks to explain. It can instead identify the presence or absence of external exit capacity and then ask whether the predicted combination with activation conditions is observed.

An inquisitorial apparatus operating inside a state able to prevent flight would illustrate high enforcement capacity joined to legally or territorially restricted exit [CITE: needs source]. A residential-school system that removed children from family relationships and language would illustrate institutional production of weak outside ties across generations [CITE: needs source]. Doctrinally strict groups situated in open societies, where members can and do leave, would provide the contrasting case of codification and active discipline without independently foreclosed exit [CITE: needs source]. These examples are illustrative hypotheses, not established historical findings in this draft. Their value lies in showing how cases could be compared on structural inputs without treating doctrine as a proxy for closure.

\subsection{Implications and testable predictions}

The framework generates several predictions that can be evaluated outside the model. First, designating an enforcement body should have no measurable effect on conduct unless that body also receives operative sanctioning privileges. The ablation supports a stronger architectural expectation: adding any single privilege to a floor with no cadre recovers no concentration, because a privilege has no designated substrate through which to operate. [VERIFY: exact single-privilege add-back estimates and the privilege count are not reported in RESULTS_TABLE.md.] Empirically, nominal offices should therefore be distinguished from offices that possess exclusive authority, reduced sanctioning costs, protection from retaliation, or material rewards.

Second, enforcement-dominant outcomes should be observed only where exit is foreclosed. Cases with highly observable codes and substantial rewards for enforcement, but open exit, should display active yet leaky enforcement. Such organizations may punish, monitor, and maintain specialized roles, but departures should prevent durable retention from following automatically. The open-boundary result gives the sharp model analogue: activation is common while capture is absent. By contrast, where independently measured exit capacity is low, otherwise similar activation conditions should be associated with retention.

Third, the mechanism is content-neutral and should apply to non-religious codified organizations. Political movements, firms, schools, residential communities, and ideological associations could exhibit the same conjunction if they make conduct observable, reward enforcement, allocate sanctioning privileges asymmetrically, and restrict outside options. This is a structural prediction, not an assertion that these types of organizations generally satisfy the conditions. Any empirical application would need case-specific evidence for each input [CITE: needs source].

Testing these predictions requires scoring cases on structural inputs blind to outcome. Researchers should code observability, enforcement reward, operative privilege, and exit capacity without using persistence, punishment concentration, or continued membership as evidence for those inputs. Outcome-blind coding would prevent the framework from becoming tautological and would preserve the distinction between conditions that activate enforcement and conditions that retain members. Comparisons should also include negative cases, especially highly codified organizations with open exit and organizations with constrained exit but weak observability or reward.

\subsection{Limitations}

The evidence comes from a single model family with a frozen specification. That discipline supports reproducibility, but it does not establish that the same surfaces will appear under alternative behavioral rules. The discrete outcome classification also depends on an activation threshold. The necessity result and dose-response remain jointly robust over the band $[0.075,0.15]$, with the canonical threshold inside that interval, but the classification should be read as a summary of continuous outcomes rather than as a threshold-free discovery.

The two dimensions are not fully orthogonal. Sealing exit raises the probability of activation in 12 of 63 cells and never lowers it; the overall active rate rises from 0.598 to 0.661. The affected cells lie at the frontier margin. Activation and retention are therefore largely separable, coupled at the frontier margin, rather than cleanly independent. This qualification does not erase the open-exit or necessity results, but it limits any claim that interventions on one dimension leave the other entirely unchanged.

Exit capacity is exogenous in the main reported analysis. An endogenous variant is reported separately as a demonstration rather than a full treatment. This choice is a safeguard against circularity, but it also means that the present model does not supply a complete theory of how legal, economic, relational, or generational processes produce closure. In addition, the reported runs contain no systematic global sensitivity analysis and no robustness analysis across network topologies. The findings are consequently conditional on the explored parameter grids and network specification.

Finally, the operational outcome classes are researcher-defined categories, not demonstrated dynamical attractors. Activation and retention thresholds organize the output into interpretable cells, but they do not prove that the simulated system has distinct long-run basins or phase-stable regimes. Claims should therefore remain at the level supported by the design: structural conditions shift enforcement activity, exit, retention, and punishment concentration in the reported model, while stronger dynamical claims await separate analysis.
